\documentclass[11pt]{amsart}
\usepackage[margin=1.1in]{geometry}
\usepackage{amsmath,amssymb,amsthm}
\usepackage{graphicx}
\usepackage[hidelinks]{hyperref}

\newtheorem{theorem}{Theorem}
\newtheorem{lemma}{Lemma}
\newtheorem{remark}{Remark}

\newcommand{\Hh}{\mathcal H}
\newcommand{\N}{\mathbb N}
\newcommand{\Tr}{\operatorname{Tr}}
\newcommand{\rank}{\operatorname{rank}}
\newcommand{\spanop}{\operatorname{span}}
\newcommand{\diag}{\operatorname{diag}}

\title[Trace-Class Triangular Projection Witness]{A Trace-Class Operator Whose Triangular Projection Is Not Trace Class}
\author{Agent lane 11}
\date{June 26, 2026}

\begin{document}
\maketitle

\section*{Source and Question}

The source paper is Alice Barbara Tumpach, \emph{Banach Poisson-Lie groups
and Bruhat-Poisson structure of the restricted Grassmannian},
arXiv:1805.03292. In Remark 1.14 on page 9, after defining the triangular
truncations \(T_-\) and \(T_{++}\), the paper recalls the classical theorem
that these truncations are unbounded on \(L_\infty(\Hh)\) and on
\(L_1(\Hh)\), but bounded on \(L_p(\Hh)\) for \(1<p<\infty\)
\cite[Remark 1.14]{Tumpach2018}. The same remark then asks for the existence
and construction of a trace-class operator whose triangular projection is not
trace class.

\begin{figure}[h]
\centering
\includegraphics[width=0.96\linewidth]{figures/open_problem_crop.png}
\caption{Remark 1.14 on page 9 of the source paper, including the
unboundedness statement and the trace-class triangular-projection question.}
\end{figure}

\section*{Proof Intuition}

Unboundedness of triangular truncation on \(S_1\) already contains the needed
witness, but not in a single line. If every finite-dimensional triangular
projection had norm bounded by one universal constant, then finite-rank
operators would give a uniformly bounded dense-domain triangular truncation,
and continuity would extend it to a bounded operator on all trace-class
operators. This contradicts the source-cited unboundedness theorem.

Thus there are finite matrix blocks whose trace norm is tiny but whose
triangular part has trace norm bounded below. Putting these blocks on disjoint
consecutive diagonal intervals gives a trace-class direct sum. Triangular
truncation acts independently on those diagonal blocks, so the first \(N\)
block compressions of the triangular part have trace norm at least \(N\). A
trace-class operator cannot have uniformly unbounded finite compressions, so
the triangular part is not trace class.

\section*{Notation}

Let \(\Hh=\ell^2(\N)\) with its usual ordered orthonormal basis
\((e_j)_{j\ge1}\). For a matrix \(A=(a_{ij})\), write
\[
  \mathcal T(A)=(b_{ij}), \qquad
  b_{ij}=
  \begin{cases}
    a_{ij}, & i\le j,\\
    0, & i>j.
  \end{cases}
\]
This is the same entrywise triangular projection as the source paper's
displayed definition of \(T_-\), up to the harmless choice of indexing order.
Write \(S_1(\Hh)\) for the trace-class operators on \(\Hh\), with norm
\(\|A\|_1=\Tr(|A|)\). If \(E\) is an \(n\)-dimensional initial coordinate
subspace, \(P_n:S_1(E)\to S_1(E)\) denotes the corresponding finite matrix
triangular projection.

\section*{Main Result}

\begin{lemma}[finite-dimensional triangular norms are unbounded]
\label{lem:finite-unbounded}
The norms \(\|P_n:S_1^n\to S_1^n\|\) are unbounded as \(n\to\infty\).
\end{lemma}

\begin{proof}
Suppose instead that \(\|P_n\|\le C\) for every \(n\). Then for every
finite-rank operator \(F\) supported on an initial coordinate subspace,
\[
  \|\mathcal T(F)\|_1\le C\|F\|_1.
\]
Since such finite-rank operators are dense in \(S_1(\Hh)\), this inequality
extends by continuity to a bounded linear map \(R:S_1(\Hh)\to S_1(\Hh)\).
For fixed \(i,j\), the coefficient functional
\(A\mapsto \langle A e_j,e_i\rangle\) is continuous on \(S_1(\Hh)\). Hence
approximating \(A\in S_1(\Hh)\) by finite-rank initial compressions shows that
the matrix coefficients of \(R A\) are exactly those of the entrywise
triangular truncation \(\mathcal T(A)\).

Thus triangular truncation would be a bounded operator on \(S_1(\Hh)\). This
contradicts the classical unboundedness theorem recalled in Remark 1.14 of
the source paper \cite{Tumpach2018}, which cites Macaev, Kwapien--Pelczynski,
Gohberg--Krein, and Arazy for this fact
\cite{Macaev1961,KwapienPelczynski1970,GohbergKrein1970,Arazy1978}.
\end{proof}

\begin{theorem}
\label{thm:witness}
There exists \(A\in S_1(\Hh)\) such that the entrywise triangular projection
\(\mathcal T(A)\) is not trace class. Moreover \(A\) can be chosen as a block
diagonal direct sum of finite-rank operators on consecutive coordinate blocks.
\end{theorem}

\begin{proof}
By Lemma~\ref{lem:finite-unbounded}, for each \(k\ge1\) there are an integer
\(n_k\) and a matrix \(C_k\in S_1^{n_k}\) such that
\[
  \|C_k\|_1=1,\qquad \|P_{n_k}C_k\|_1\ge 2^k.
\]
Set \(B_k=2^{-k}C_k\). Then
\[
  \|B_k\|_1=2^{-k}, \qquad \|P_{n_k}B_k\|_1\ge 1.
\]

Choose pairwise disjoint consecutive intervals
\[
  I_k=\{m_k+1,\ldots,m_k+n_k\}\subset\N
\]
and order-preserving unitaries \(U_k:\mathbb C^{n_k}\to \ell^2(I_k)\). Define
\[
  A_k=U_kB_kU_k^*,\qquad
  A=\bigoplus_{k=1}^{\infty} A_k
\]
with respect to the orthogonal decomposition
\(\Hh=\bigoplus_k \ell^2(I_k)\) plus any unused coordinates. Since
\[
  \sum_{k=1}^\infty \|A_k\|_1
  =\sum_{k=1}^\infty 2^{-k}<\infty,
\]
the block diagonal operator \(A\) is trace class.

Because the blocks are consecutive and all off-block matrix coefficients of
\(A\) vanish, triangular truncation preserves the block diagonal form:
on the block \(I_k\), the compression of \(\mathcal T(A)\) is
\[
  U_k(P_{n_k}B_k)U_k^*.
\]
Let \(Q_N\) be the orthogonal projection onto the sum of the first \(N\)
blocks. Then
\[
\begin{aligned}
  \|Q_N\mathcal T(A)Q_N\|_1
    &= \sum_{k=1}^{N}\|U_k(P_{n_k}B_k)U_k^*\|_1  \\
    &= \sum_{k=1}^{N}\|P_{n_k}B_k\|_1
     \ge N .
\end{aligned}
\]
If \(\mathcal T(A)\) were trace class, then
\(\|Q_N\mathcal T(A)Q_N\|_1\le \|\mathcal T(A)\|_1\) for every \(N\), an
impossibility. Therefore \(\mathcal T(A)\notin S_1(\Hh)\).
\end{proof}

\section*{Consequence for the Source Question}

Theorem~\ref{thm:witness} supplies the trace-class operator requested in
Remark 1.14. The construction is canonical once one selects a sequence of
finite-dimensional matrices witnessing the unboundedness of triangular
projection on \(S_1^n\). The same paper later uses unboundedness of triangular
truncation to choose a sequence \(K_n\in L_1\) with bounded trace norm and
triangular truncations of arbitrarily large trace norm
\cite[Proposition 7.1 proof]{Tumpach2018}; the present argument upgrades this
``unbounded sequence'' phenomenon to one single trace-class operator by the
block-direct-sum device.

\section*{Verification Notes}

No numerical verification is needed for the proof. The two functional-analytic
checks are:
\begin{enumerate}
\item finite-dimensional triangular norms must be unbounded, otherwise the
dense finite-rank truncation extends boundedly to all \(S_1\);
\item for a trace-class operator \(C\), finite-rank compressions satisfy
\(\|Q C Q\|_1\le \|C\|_1\), while the constructed triangular truncation has
compressions with norms tending to infinity.
\end{enumerate}
The script in \texttt{code/render\_open\_problem\_crop.py} only regenerates
the source-page crop.

\section*{Novelty Check}

Before promotion, the run checked the lightweight registry, solution, attempt,
and proof-gap indexes for arXiv:1805.03292 and the phrases
``triangular projection'', ``triangular truncation'', ``trace class operator'',
``not trace class'', ``Macaev'', ``Kwapien'', ``Pelczynski'', ``Gohberg'',
``Krein'', and ``Arazy''. No duplicate packet was found. A local parsed-corpus
search found the source question, the same paper's later use of an unbounded
sequence of trace-class witnesses, and related later arXiv records on
triangular projection for \(0<p<1\), but no separate answer to the
\(S_1\) single-operator question. Web searches for exact phrases including
``trace class operator triangular projection not trace class'', ``triangular
truncation trace class not trace class'', and ``Tumpach triangular projection
trace class'' did not reveal a later paper explicitly answering Remark 1.14.

\section*{Limitations}

This packet answers only the trace-class triangular-projection question in
Remark 1.14. It does not address the separate broad problem, also present in
the source paper, about integrating Banach Lie bialgebras into Banach
Poisson-Lie groups. The construction is explicit at the operator-assembly
level but depends on choosing finite-dimensional matrices from the classical
unboundedness theorem for triangular projections on \(S_1^n\).

\section*{Human-Review Recommendation}

Review as a likely valid full answer to the trace-class triangular-projection
open problem in Remark 1.14. The main point to verify is the passage from
source-cited unboundedness on \(S_1(\Hh)\) to unbounded finite-dimensional
triangular projection norms in Lemma~\ref{lem:finite-unbounded}.

\begin{thebibliography}{9}

\bibitem{Tumpach2018}
A. B. Tumpach,
\emph{Banach Poisson-Lie groups and Bruhat-Poisson structure of the
restricted Grassmannian},
arXiv:1805.03292; Comm. Math. Phys. 373 (2020), 795--858.

\bibitem{Arazy1978}
J. Arazy,
\emph{Some remarks on interpolation theorems and the boundness of the
triangular projection in unitary matrix spaces},
Integral Equations Operator Theory 1 (1978), 453--495.

\bibitem{GohbergKrein1970}
I. C. Gohberg and M. G. Krein,
\emph{Theory and Applications of Volterra Operators on Hilbert Space},
Translations of Mathematical Monographs, vol. 24, American Mathematical
Society, 1970.

\bibitem{KwapienPelczynski1970}
S. Kwapien and A. Pelczynski,
\emph{The main triangle projection in matrix spaces and its applications},
Studia Math. 34 (1970), 43--68.

\bibitem{Macaev1961}
V. I. Macaev,
\emph{Volterra operators produced by perturbation of self-adjoint operators},
Soviet Math. Dokl. 2 (1961), 1013--1016.

\end{thebibliography}

\end{document}
